\documentclass[11pt,a4paper]{article}
\usepackage[T2A]{fontenc}
\usepackage[utf8]{inputenc}
\usepackage[russian]{babel}
\usepackage{amsmath, amssymb, amsfonts, amsthm}
\usepackage{geometry}
\geometry{top=2.0cm, bottom=2.0cm, left=2.0cm, right=2.0cm}
\usepackage{hyperref}
\usepackage{booktabs}
\usepackage{tabularx}
\hypersetup{colorlinks=true, linkcolor=blue, citecolor=blue, urlcolor=blue}

\newtheorem{theorem}{Теорема}
\newtheorem{lemma}[theorem]{Лемма}
\theoremstyle{definition}
\newtheorem{definition}{Определение}
\theoremstyle{remark}
\newtheorem{remark}{Замечание}

\begin{document}

\title{\textbf{\Large Инженерная модель структуры нейтрона:}\\
\large электромагнитное расщепление $\Delta M_{np}$, акустический эффект Парселла и топологический предел странности\\[2mm]
\normalsize \textit{Серия: Расчетные методы топологической эластодинамики континуума. Статья 03}}

\author{\textbf{Дмитрий Михайленко}\\[1mm]
\small Исследовательская группа топологической физики континуума}

\date{\today}

\maketitle

\begin{abstract}
\noindent В рамках модели упругого 3D-континуума представлена замкнутая расчетная схема статических и динамических характеристик нейтрона. Нейтрон описывается как базовый трилистник $T(3,2)$, локализованный внутри сферической фазовой доменной стенки $S^2$. Положительный знак разности масс нуклонов ($M_n > M_p$) выведен из механической работы упругого сжатия керна против модуля сдвига $G_0$. На основе нуклонного форм-фактора $f = R_m / r_c = 3/4$ выведено аналитическое соотношение для разности масс $\Delta M_{np} = \frac{3}{16}\alpha M_p (1 + \alpha/\pi) \approx 1.2868\text{ МэВ}$, согласующееся с экспериментом до $0.51\%$ без феноменологических масс кварков. Экспериментальное расхождение времени жизни свободного нейтрона между пучковым методом ($\tau_{\mathrm{beam}} \approx 887.7\text{ с}$) и гравитационными ловушками ($\tau_{\mathrm{bottle}} \approx 878.4\text{ с}$) разрешается через акустический эффект Парселла: макроскопическая граница вакуумной полости ловушки модифицирует фазовый объем испускаемых волн деформации, давая сдвиг $\Delta \tau = \tau_{\mathrm{beam}} \cdot (\frac{4}{3}\alpha) \approx 8.64\text{ с}$ (расчетное значение в ловушке: $879.06\text{ с}$, расхождение с UCNτ $+0.66\text{ с}$). Топологическое число лепестков трилистника $p=3$ строго ограничивает максимальную странность барионов величиной $|S| \le 3$, а масса $\Lambda^0$-гиперона воспроизводится с точностью $0.12\%$. Все математические соотношения верифицированы в интерактивном прувере Lean 4 без допущений (\texttt{sorry}).
\end{abstract}

\vspace{2mm}
\noindent\textbf{Ключевые слова:} нейтрон, разность масс $\Delta M_{np}$, аномалия времени жизни, эффект Парселла, ультрахолодные нейтроны, предел странности, Lean 4.

\section{Введение: статус нейтрона в модели континуума}
В Стандартной модели разность масс нейтрона и протона $\Delta M_{np} = M_n - M_p \approx 1.29333\text{ МэВ}$ является тонким балансом между электромагнитным самодействием (которое стремится сделать протон тяжелее) и разностью токовых масс кварков ($m_d - m_u \approx 2.5\text{ МэВ}$), вводимой исключительно вручную. 

Кроме того, более двух десятилетий сохраняется кризис времени жизни нейтрона:
\begin{itemize}
    \item Пучковый метод (NIST): регистрация протонов распада в открытом пучке дает $\tau_{\mathrm{beam}} = 887.7 \pm 1.2\text{ с}$;
    \item Метод ультрахолодных нейтронов (UCN, $\mathrm{UCN}\tau$ LANL): подсчет выживших нейтронов в магнитно-гравитационных и материальных ловушках дает $\tau_{\mathrm{bottle}} = 878.4 \pm 0.5\text{ с}$ (в новейших сериях $877.75 \pm 0.36\text{ с}$).
\end{itemize}
Расхождение $\Delta\tau \approx 8.5\dots 9.3\text{ с}$ превышает $4.5\sigma$, порождая спекуляции о распадах в зеркальную или темную материю.

Настоящая статья решает обе задачи из геометрических свойств упругого вакуума с использованием ранее доказанных инвариантов трилистника $T(3,2)$ (Статьи 01 и 02).

\section{Модель нейтрона: трилистник под доменной стенкой $S^2$}
В открытом пространстве трилистник $T(3,2)$ обладает свободными радиальными границами, формируя электромагнитное поле свободного заряда $+1e$ (протон). 

Нейтрон возникает, когда топологический заряд экранируется замкнутой сферической доменной стенкой $S^2$ с вихревой поляризацией обратного знака:
\begin{equation}
    \oint_{S^2} \nabla \theta \cdot \mathrm{d}\mathbf{S} = 0 \implies Q_{\mathrm{total}} = 0.
\end{equation}

\begin{theorem}[Механический знак расщепления $\Delta M_{np} > 0$]
Наличие внешней сферической границы $S^2$ требует приложения радиального сжимающего напряжения $\sigma_{rr} < 0$. В упругом квазинесжимаемом континууме объемная работа сжатия против модуля сдвига $G_0$ положительна:
\begin{equation}
    \Delta E_{\mathrm{elast}} = \frac{1}{2} \int_{V} \frac{\sigma_{rr}^2}{G_0} \, \mathrm{d}V > 0 \implies M_n > M_p.
\end{equation}
\end{theorem}
В отличие от феноменологии кварков, где знак $m_d > m_u$ постулируется, в механике континуума нейтрон обязан быть тяжелее протона, поскольку он представляет собой напряженное (сжатое) состояние того же узла.

\section{Дедуктивный вывод разности масс $\Delta M_{np}$}
Из Статьи 02 известно, что отношение механического радиуса ядра $R_m = 3\lambda_p$ к электромагнитному зарядовому радиусу $r_c = 4\lambda_p$ строго равно форм-фактору $f = 3/4$.

\begin{theorem}[Электромагнитно-упругое расщепление масс]
\label{thm:mass_split}
Разность масс между сжатым состоянием (нейтрон) и свободным состоянием (протон) определяется геометрической проекцией волнового импеданса $\alpha$ на квадрупольную деформацию четырех спинорных степеней свободы:
\begin{equation}
    \Delta M_{np} = C_{np} \, \alpha M_p \left( 1 + \frac{\alpha}{\pi} \right), \quad C_{np} = f \times \frac{1}{4} = \frac{3}{4} \times \frac{1}{4} = \frac{3}{16}.
\end{equation}
\end{theorem}
\begin{proof}
1. Энергия экранирования сосредоточена в переходном сферическом слое между $R_m$ и $r_c$. Доля упругой энергии ядра, доступная для поляризации, задается геометрическим форм-фактором $f = R_m / r_c = 3/4$.
2. Экранирование дипольного момента в континууме с $\mathrm{Spin}(3) \cong \mathrm{SU}(2)$ распределяется по $N_{\mathrm{spin}} = 4$ вещественным спинорным компонентам матрицы плотности, что дает делитель $1/4$.
3. Безразмерный коэффициент связи равен $C_{np} = \frac{3}{4} \cdot \frac{1}{4} = \frac{3}{16}$.
4. Сдвиг границы раздела фаз генерирует стандартную акустическую радиационную поправку первого порядка по импедансу границы $\delta_{\mathrm{rad}} = \alpha/\pi$.
5. Итоговая разность масс равна $\Delta M_{np} = \frac{3}{16}\alpha M_p (1 + \alpha/\pi)$.
\end{proof}

\subsection{Численный расчет}
Используя базовые параметры $M_p = 938.272088\text{ МэВ}$ и $\alpha^{-1} = 137.035999177$:
\begin{equation}
    \Delta M_{np}^{(0)} = \frac{3}{16} \times \frac{938.272088}{137.035999177} = 1.283794\text{ МэВ}.
\end{equation}
С учетом пограничной поправки $(1 + \alpha/\pi) \approx 1.002323$:
\begin{equation}
    \Delta M_{np} = 1.283794 \times 1.002323 = \mathbf{1.28678\text{ МэВ}}.
\end{equation}
Экспериментальное значение CODATA 2022: $\Delta M_{np}^{\mathrm{exp}} = 1.293332\text{ МэВ}$. 
Относительная ошибка составляет $-0.51\%$ без привлечения феноменологических параметров.

\section{Аномалия времени жизни: акустический эффект Парселла}
Распад свободного нейтрона $n \to p + e^- + \bar{\nu}_e$ сопровождается излучением солитонной петли (электрона) и фазовой 1D-нити (нейтрино). Скорость бета-распада по золотому правилу Ферми пропорциональна фазовому объему излучаемых мод континуума: $\Gamma \propto \rho(\omega)$.

В 1946 году Э. Парселл показал, что спонтанное излучение источника изменяется при помещении его в резонатор с отражающими границами. В экспериментах с ультрахолодными нейтронами (UCN) частицы удерживаются внутри макроскопической вакуумной полости; внешняя граница ловушки создает резкий скачок акустико-электромагнитного импеданса.

\begin{theorem}[Сдвиг времени жизни Парселла в ловушках UCN]
Наличие внешней отражающей границы ловушки сжимает фазовый объем для длинноволновых мод деформации. Относительное изменение скорости распада масштабируется обратным форм-фактором нуклона $1/f = 4/3$ и волновым импедансом $\alpha$:
\begin{equation}
    \frac{\Delta \Gamma}{\Gamma} = \frac{\Gamma_{\mathrm{bottle}} - \Gamma_{\mathrm{beam}}}{\Gamma_{\mathrm{beam}}} = \frac{4}{3}\alpha.
\end{equation}
Соответствующий сдвиг времени жизни равен:
\begin{equation}
    \Delta \tau = \tau_{\mathrm{beam}} - \tau_{\mathrm{bottle}} = \tau_{\mathrm{beam}} \cdot \left( \frac{4}{3}\alpha \right).
\end{equation}
\end{theorem}

\subsection{Численная верификация}
Используя пучковое значение NIST $\tau_{\mathrm{beam}} = 887.7 \pm 1.2\text{ с}$:
\begin{equation}
    \Delta \tau = 887.7\text{ с} \times \left( \frac{4}{3 \times 137.035999177} \right) = 887.7 \times 0.0097298 = \mathbf{8.637\text{ с}}.
\end{equation}
Расчетное время жизни нейтрона в закрытой ловушке:
\begin{equation}
    \tau_{\mathrm{bottle}}^{\mathrm{theor}} = 887.7 - 8.64 = \mathbf{879.06\text{ с}}.
\end{equation}
Экспериментальные данные гравитационных ловушек:
\begin{itemize}
    \item PDG среднее по ловушкам: $\tau_{\mathrm{bottle}}^{\mathrm{exp}} = 878.4 \pm 0.5\text{ с}$ (расхождение с моделью: $0.66\text{ с}$);
    \item Эксперимент $\mathrm{UCN}\tau$ (LANL): $877.75 \pm 0.36\text{ с}$ (расхождение: $1.3\text{ с}$).
\end{itemize}
Таким образом, «аномалия времени жизни нейтрона» объясняется классическим эффектом граничных условий волновода без привлечения гипотез о темных распадах.

\section{Топологический предел странности и гиперон $\Lambda^0$}
В кварковой модели странность $|S|$ вводится как квантовое число $s$-кварка. В модели ТЭВ странность представляет собой вторичное внутреннее закручивание пучности трилистника на мюонную частотную моду $\omega_\mu$.

\begin{theorem}[Предел странности барионов]
Поскольку базовый узел $T(3,2)$ имеет ровно $p=3$ полоидальные пучности, максимальное число одновременно возбужденных странных мод строго ограничено:
\begin{equation}
    |S|_{\mathrm{max}} = p = 3.
\end{equation}
\end{theorem}
Это чисто геометрически объясняет, почему в природе существуют барионы со странностью $S = 0, -1, -2, -3$ ($\Omega^-$), но фундаментально не существует барионов со странностью $|S| \ge 4$.

\subsection{Масса легчайшего гиперона $\Lambda^0$}
Одиночное возбуждение странности ($S = -1$) соответствует вовлечению массы мюона $m_\mu$ с учетом фазового угла связи пучности $1 + 2/3 = 5/3$:
\begin{equation}
    M_{\Lambda^0} = M_p + \frac{5}{3} m_\mu.
\end{equation}
Подставляя $M_p = 938.272\text{ МэВ}$ и $m_\mu = 105.658\text{ МэВ}$:
\begin{equation}
    M_{\Lambda^0}^{\mathrm{theor}} = 938.272 + \frac{5}{3}(105.658) = 938.272 + 176.097 = \mathbf{1114.37\text{ МэВ}}.
\end{equation}
Эксперимент PDG: $M_{\Lambda^0}^{\mathrm{exp}} = 1115.683 \pm 0.006\text{ МэВ}$. Точность предсказания: $0.12\%$ ($-1.31\text{ МэВ}$).

\section{Сводные результаты}

\begin{table}[h]
\centering
\small
\caption{Сводка результатов Статьи 03 по параметрам нейтрона и гиперона}
\label{tab:neutron_summary}
\vspace{2mm}
\begin{tabular}{lcccc}
\toprule
\textbf{Величина} & \textbf{Формула ТЭВ} & \textbf{Модель} & \textbf{Эксперимент} & \textbf{Невязка} \\
\midrule
Разность масс $\Delta M_{np}$ & $\frac{3}{16}\alpha M_p (1 + \alpha/\pi)$ & \textbf{1.2868 МэВ} & $1.293332$ МэВ & $-0.51\%$ \\
Сдвиг Парселла $\Delta\tau$ & $\tau_{\mathrm{beam}} \cdot (\frac{4}{3}\alpha)$ & \textbf{8.64 с} & $8.5\dots 9.3$ с & В пределах ошибок \\
Время жизни в ловушке & $\tau_{\mathrm{beam}} - \Delta\tau$ & \textbf{879.06 с} & $878.4 \pm 0.5$ с & $0.07\%$ \\
Максимальная странность & $|S| \le p$ & \textbf{3} & $\Omega^-$ ($S=-3$) & Тождественно \\
Масса $\Lambda^0$-гиперона & $M_p + \frac{5}{3}m_\mu$ & \textbf{1114.37 МэВ} & $1115.683$ МэВ & $-0.12\%$ \\
\bottomrule
\end{tabular}
\end{table}

\section{Заключение}
Модель показывает, что ключевые свойства нейтрона — устойчивое превышение массы над протоном, величина расщепления $\Delta M_{np}$, расхождение времени жизни в ловушках UCN и порог странности — являются прямыми следствиями гидродинамики и топологии сжатого трилистника $T(3,2)$.

\begin{thebibliography}{99}

\bibitem{Purcell1946}
E.~M.~Purcell, \textit{Spontaneous emission probabilities at radio frequencies}, Phys. Rev. \textbf{69}, 681 (1946).

\bibitem{NIST2013}
A.~T.~Yue, M.~S.~Dewey, D.~M.~Gilliam, J.~S.~Nico, F.~E.~Wietfeldt, C.~A.~Fehsenfeld, and W.~M.~Snow, \textit{Improved determination of the neutron lifetime}, Phys. Rev. Lett. \textbf{111}, 222501 (2013).

\bibitem{UCNtau2021}
F.~M.~Gonzalez \textit{et al.} (UCN$\tau$ Collaboration), \textit{Improved neutron lifetime measurement with UCN$\tau$}, Phys. Rev. Lett. \textbf{127}, 162501 (2021).

\bibitem{Serebrov2008}
A.~P.~Serebrov \textit{et al.}, \textit{Neutron lifetime measurements with a large gravitational trap for ultracold neutrons}, Phys. Rev. C \textbf{78}, 035505 (2008).

\bibitem{Czarnecki2018}
A.~Czarnecki, W.~J.~Marciano, and A.~Sirlin, \textit{Neutron lifetime and axial coupling connection}, Phys. Rev. Lett. \textbf{120}, 202002 (2018).

\bibitem{GellMann1962}
M.~Gell-Mann, \textit{Symmetries of baryons and mesons}, Phys. Rev. \textbf{125}, 1067--1084 (1962).

\bibitem{PDG2024}
S.~Navas \textit{et al.} (Particle Data Group), \textit{Review of Particle Physics}, Phys. Rev. D \textbf{110}, 030001 (2024).

\end{thebibliography}

\end{document}
